% 中文摘要

超高真空是现代高端科学实验与先进制造技术的关键基础条件，其测量精度直接关系到相关实验结果与工艺过程的可靠性。传统电离规类真空测量方法依赖气体电离过程，易受气体组分、表面效应及器件老化等因素影响，存在固有系统误差及长期稳定性不足等问题，难以满足高精度、可溯源真空计量的需求。基于冷原子碰撞损失机制的量子真空测量技术，通过利用背景气体与冷原子之间的碰撞过程，将原子损失率与真空压强建立直接对应关系，其测量结果可溯源至基本物理常数，因而具有高精度与良好长期稳定性的优势，近年来逐渐成为真空计量领域的重要研究方向。然而，现有研究主要集中于测量原理验证与计量模型构建，在装置小型化、系统工程化集成及实际应用适配方面仍面临显著挑战，制约了该技术的进一步推广。

本文围绕小型化冷原子真空计量系统的设计与实验验证开展研究。针对传统冷原子实验装置体积大、结构复杂的问题，对系统进行了整体优化设计，包括紧凑型光学系统、永磁体—电磁线圈复合磁场结构以及小型化超高真空腔体。基于铷原子实验平台，构建了完整的冷原子实验装置并实现磁光阱装载。在此基础上，将开展磁阱中原子寿命测量相关实验研究，通过记录原子数随时间的衰减过程提取损失率参数，并建立其与背景气体压强之间的定量关系。同时，对系统测量结果的一致性、重复性及主要误差来源进行了系统分析。

实验结果表明，所设计的小型化冷原子真空计量系统能够稳定实现磁光阱装载，并表现出良好的测量一致性与重复性。相关研究为冷原子真空计量装置的小型化实现提供了实验依据与技术参考，对推动量子真空计量技术向工程化与实际应用发展具有积极意义。


